\section{Native REV\TeX{} transcription of current Section~4 source block (pp.~27--67)}
This appendix begins the native rewrite of the March~16,~2026 source manuscript by replacing the Section~4 archival facsimile pages with a native theorem-facing transcription. The focus of this phase is the flowed continuum state and the dyadic-RG / no-subsequence convergence infrastructure culminating in source Theorem~4.92 and Corollaries~4.93--4.95.

\subsection{Group-uniform A1 input and the dyadic GF step-scaling lane}
The source Section~4 begins by isolating the one-plaquette coefficient input and the dyadic gradient-flow step-scaling structure used throughout Lane~C.

\begin{theorem}[Source Theorem~4.2: group-uniform A1 input used on the primary spine]
For every compact connected simple gauge group $G$ and every fixed finite-dimensional unitary Wilson representation $\rho$, the one-plaquette coefficients satisfy the crossover bound of the compact-simple A1 theorem. Equivalently, on the primary Route~1 spine every use of the A1 gate is available with group-dependent constants $(c_{A1}(G,\rho),\beta_{A1}(G,\rho))$ and no remaining Lie-type restriction.
\end{theorem}

\begin{proposition}[Source Proposition~4.6: one-loop coefficient is scheme invariant]
Let $u$ and $u'$ be two renormalized couplings related near $0$ by an analytic reparametrization
\[
u' = u + a_1 u^2 + a_2 u^3 + O(u^4).
\]
If
\[
\beta(u) = -2b_0u^2 + O(u^3),
\]
then
\[
\beta'(u') = -2b_0 (u')^2 + O((u')^3).
\]
In particular, the coefficient $b_0$ is invariant under the reparametrization.
\end{proposition}

\begin{lemma}[Source Lemma~4.7: beta-remainder implies step-scaling remainder]
Assume there exist constants $u_0>0$ and $B\ge 0$ such that
\[
|\beta_{\mathrm{GF}}(u)+2b_0u^2| \le B u^3
\qquad (0\le u\le u_0),
\]
and assume $u_0$ is small enough that $D:=1-2b_0u_0\log 2>0$. Then for every $u\in[0,u_0]$, the factor-$2$ step-scaling map satisfies
\[
\Sigma(u)=u+2b_0u^2\log 2 + R_\Sigma(u),
\qquad
|R_\Sigma(u)|\le C_\Sigma u^3
\]
with explicit constant
\[
C_\Sigma:= \frac{(2b_0\log 2)^2}{D} + \frac{B\log 2}{D^3}.
\]
\end{lemma}

\subsection{One-step analytic RG map and dyadic beta remainder}
The middle of the source Section~4 turns the Lane~C dyadic RG into a scale-uniform analytic map with explicit Taylor remainder bounds.

\begin{proposition}[Source Proposition~4.9: analytic RG map with explicit Taylor remainder bounds]
There exist constants
\[
r_0>0,\qquad \kappa\in(0,1),\qquad A_2\in\mathbb R,\qquad C_u>0,\qquad C_K>0,
\]
and linear operators $L_j:\mathcal K_j\to \mathcal K_{j+1}$ with $\|L_j\|_{\mathrm{op}}\le \kappa$, such that for all $j$ and all $(u,K)\in B_j(r_0)$ the RG step satisfies
\begin{align}
u' &= u + A_2 u^2 + \delta_u(u,K),\\
K' &= L_j K + \delta_K(u,K),
\end{align}
with bounds
\begin{align}
|\delta_u(u,K)| &\le C_u \bigl(|u|^3 + |u|\,\|K\|_j\bigr),\\
\|\delta_K(u,K)\|_{j+1} &\le C_K \bigl(|u|^2 + \|K\|_j^2\bigr).
\end{align}
All constants are uniform in the dyadic scale.
\end{proposition}

\begin{proposition}[Source Proposition~4.10: entrance estimate]
There exists an entrance scale $j_\ast$ and a constant $D>0$ such that
\[
\|K_{j_\ast}\|_{j_\ast}\le D u_{j_\ast}^2,
\qquad
|u_{j_\ast}|\le r_0/2.
\]
\end{proposition}

\begin{proposition}[Source Proposition~4.11: quadratic coefficient identification]
The RG marginal coefficient satisfies
\[
A_2 = 2b_0 \log 2,
\qquad b_0>0
\]
for pure Yang--Mills. In the load-bearing chain only the positivity $b_0>0$ is used downstream.
\end{proposition}

\begin{lemma}[Source Lemma~4.13: irrelevant sector stays $O(u^2)$]
Assume the analytic RG map above. Define the cone region
\[
\mathcal C_j(M):=\{(u,K)\in B_j(r_0): \|K\|_j\le M u^2\}.
\]
If
\[
M\ge \frac{2C_K}{1-\kappa}
\]
and $r_0$ is sufficiently small, then
\[
(u,K)\in \mathcal C_j(M)\implies (u',K')\in \mathcal C_{j+1}(M).
\]
\end{lemma}

\begin{corollary}[Source Corollary~4.14: uniform irrelevant bound along the dyadic trajectory]
With
\[
M:=\max\!\left\{D,\frac{2C_K}{1-\kappa}\right\},
\]
one has for all $j\ge j_\ast$,
\[
\|K_j\|_j \le M u_j^2
\]
provided the RG stays in the small-coupling ball.
\end{corollary}

\begin{theorem}[Source Theorem~4.15: dyadic step remainder bound in RG seminorm topology]
Assume the preceding RG hypotheses. Then for all $j\ge j_\ast$ and all sufficiently small $u_j$,
\[
|u_{j+1}-u_j-A_2u_j^2|\le C_\Sigma u_j^3,
\qquad
C_\Sigma:=C_u(1+M).
\]
The constant $C_\Sigma$ is uniform in $j$ and hence uniform in $a/L_j$.
\end{theorem}

\begin{corollary}[Source Corollary~4.16: uniform dyadic beta-remainder bound]
Define the dyadic beta function by
\[
\beta^{(2)}(u_j):=-\frac{u_{j+1}-u_j}{\log 2}.
\]
Then for all $j\ge j_\ast$ and all sufficiently small $u_j$,
\[
|\beta^{(2)}(u_j)+2b_0u_j^2| \le Bu_j^3,
\qquad
B:=\frac{C_\Sigma}{\log 2}
= \frac{C_u(1+M)}{\log 2},
\]
uniformly in the dyadic scale.
\end{corollary}

\subsection{Uniform convexity and the extracted-irrelevant scaling gap}
The source Section~4 next reduces the one-step fluctuation estimates to finite-dimensional uniform convexity on the good-chart domain together with the canonical extracted-irrelevant scaling gap.

\begin{lemma}[Source Lemma~4.40: explicit lower bound for the tree-gauge curl coercivity]
For the $2\times 2\times 2\times 2$ coarse block with the canonical tree gauge, the reduced integer matrix
\[
Q:=D_{\mathrm{red}}^\ast D_{\mathrm{red}}
\]
has smallest eigenvalue $\lambda_\ast$ satisfying
\[
\lambda_\ast > \frac14.
\]
Consequently the same lower bound holds for the full $G$ operator $D_{\mathrm{red}}^\ast D_{\mathrm{red}}$, since it is block-diagonal over Lie components.
\end{lemma}

\begin{theorem}[Source Theorem~4.44: uniform convexity on the chart domain]
Fix $\delta\le \min\{\pi/4,\delta_{\mathrm{Haar}}(G)\}$ and define
\[
r_\ast := \sup_{\|A\|_\infty\le \delta}\max_{p\in P(Q)} \|F_p(A)\|.
\]
Assume the BCH derivative bounds and the smallness conditions
\[
\varepsilon_1(\delta)\le \frac12\sqrt{\lambda_\ast},
\qquad
M(r_\ast;G)\varepsilon_2(\delta)r_\ast \le \frac{\alpha_{\min}}{8}m(r_\ast;G)\lambda_\ast.
\]
Then for all $A$ with $\|A_e\|\le \delta$,
\[
\nabla^2 H_{j+1}(A)\succeq \varrho I,
\qquad
\varrho := \alpha_{\min}\beta\,\frac{m(r_\ast;G)\lambda_\ast}{8}.
\]
In particular, $\varrho>0$ is uniform in scale $j$ provided $\beta\ge \beta_{\min}>0$ and $\delta$ is fixed.
\end{theorem}

\begin{lemma}[Source Lemma~4.46: scaling gap after extraction]
Assume the extraction removes all gauge-invariant local components of engineering dimension $\le 4$. Then every remaining local monomial has dimension at least $6$, so the minimal dyadic scaling gap is
\[
\Delta_{\min}=6-4=2.
\]
Under deterministic dyadic rescaling one therefore has
\[
\|\mathrm{Rescale}(K)\|_{j+1}\le 2^{-\Delta_{\min}}\|K\|_j = \frac14\|K\|_j
\]
for extracted irrelevants, provided the seminorm is defined with canonical scaling weights.
\end{lemma}

\begin{theorem}[Source Theorem~4.66: uniform dyadic beta-remainder bound for the patched RG]
Assume:
\begin{enumerate}
\item the dyadic RG step is the chart-truncated map $R_j^{\mathrm{good}}$,
\item extraction removes all relevant/marginal components up to the chosen depth,
\item the identification $A_2=2b_0\log 2$ holds with $b_0>0$,
\item the project seminorm axioms hold with scale-independent constants.
\end{enumerate}
Then there exist $u_0,B>0$, independent of the dyadic scale $j$, such that for all scales in the asymptotically free window and all $u\in[0,u_0]$,
\[
|\beta^{(2)}(u)+2b_0u^2| \le Bu^3,
\]
with explicit choice
\[
B = \frac{C_u(1+M)}{\log 2},
\qquad
M:=\max\!\left\{D,\frac{2C_K}{1-\kappa}\right\}.
\]
\end{theorem}

\subsection{Stable manifold, injective tuning, and observable Lipschitz control}
The last large block of the source Section~4 constructs the UV-stable irrelevant trajectory, proves injective tuning of the renormalization map at fixed physical scale, and establishes the observation-scale Lipschitz dependence needed for uniqueness of the continuum state.

\begin{theorem}[Source Theorem~4.72: stable manifold for the patched RG]
Let $(u_j)_{j\le 0}$ be an admissible asymptotically free trajectory. Then:
\begin{enumerate}
\item for each fixed $j\le 0$, the limit
\[
K_j^{(\infty)} := \lim_{N\to \infty,\,N\ge |j|} K_j^{(N)}
\]
exists in $\mathcal K_j$ and satisfies the RG recursion
\[
K_{j+1}^{(\infty)} = \mathcal K_j(u_j,K_j^{(\infty)});
\]
\item one has $K_j^{(\infty)}\to 0$ in norm as $j\to -\infty$;
\item the UV-stable irrelevant trajectory is unique;
\item the limit is universal with respect to bounded UV initialization in the irrelevant sector.
\end{enumerate}
Moreover,
\[
\|K_j^{(\infty)}\|_j \le C_F \sum_{m=-\infty}^{j-1}\kappa^{j-1-m}u_m^2.
\]
\end{theorem}

\begin{theorem}[Source Theorem~4.78: injectivity/monotonicity of the renormalization map]
Assume the one-step analytic remainder bounds for the step-scaling maps. Then there exists a uniform asymptotically free interval $[0,u^\circ]$ such that for every cutoff depth $N$ the renormalization map
\[
\rho_N(u_{-N}) := u_0
\]
is $C^1$ and strictly increasing on $[0,u^\circ]$. In particular, for every target $u_\mu$ in the image of $\rho_N$ there exists a unique tuned UV coordinate $u_{-N}$ with $\rho_N(u_{-N})=u_\mu$.
\end{theorem}

\begin{theorem}[Source Theorem~4.86: bulk observable Lipschitz in the RG data]
Fix a finite observable set $\mathcal F\subset A_{\mathrm{flow}}$ and choose the observation scale $j_{\mathrm{obs}}$ attached to $\mathcal F$. Then for any two datasets $(u,K),(\tilde u,\tilde K)$ in the observation-scale RG ball,
\[
\max_{F\in \mathcal F}|G_F(u,K)-G_F(\tilde u,\tilde K)|
\le L_{\mathcal F}\bigl(|u-\tilde u|+\|K-\tilde K\|_{j_{\mathrm{obs}}}\bigr),
\]
with explicit constant $L_{\mathcal F}$ independent of the lattice spacing in the patched construction.
\end{theorem}

\subsection{No-subsequence convergence and the flowed-state packet}
The source Section~4 ends by assembling the no-subsequence convergence theorem for the patched outputs and its identification with the physical Wilson continuum state.

\begin{theorem}[Source Theorem~4.92: no-subsequence convergence for the flowed continuum state]
Assume the uniform boundedness/precompactness package, injective tuning, stable-manifold uniqueness, and observation-scale Lipschitz control. Fix a target renormalized coupling value in the asymptotically free interval, and for each lattice spacing tune the UV coupling so that the renormalization condition holds. Then:
\begin{enumerate}
\item for every finite set $\mathcal F\subset A_{\mathrm{flow}}$, the joint values $\{\omega_a^{\mathrm{RG}}(F)\}_{F\in\mathcal F}$ converge as $a\to 0$ without subsequence extraction;
\item there exists a unique state $\omega$ on $A_{\mathrm{flow}}$ such that
\[
\omega(F)=\lim_{a\to 0}\omega_a(F)=\lim_{a\to 0}\omega_a^{\mathrm{RG}}(F)
\qquad (F\in A_{\mathrm{flow}});
\]
\item the limit state is Euclidean invariant and reflection invariant, and inherits the algebraic identities preserved under pointwise limits.
\end{enumerate}
\end{theorem}

\begin{corollary}[Source Corollary~4.93: unique continuum Euclidean state on $\mathcal A_{\mathrm{flow}}$]
Under the hypotheses of Theorem~\ref{thm:C3U}, the tuned unpatched Wilson expectations $\omega_a$ converge (without subsequences) to a unique Euclidean-invariant state $\omega$ on $\mathcal A_{\mathrm{flow}}$ for all flow times $t\in(0,t_\ast]$.
\end{corollary}

\begin{proof}
This is immediate from Theorem~\ref{thm:C3U}(b)--(c).
\end{proof}

\begin{corollary}[Source Corollary~4.94: tuned dyadic Wilson limit on the bounded positive-time algebra]
Fix a target value $u_{\mathrm{ren}}\in (0,u_\ast]$ and the associated tuned dyadic trajectory
\[
a_n=\ell_0 2^{-n},
\qquad
\beta_n = \beta(a_n;u_{\mathrm{ren}}).
\]
Then for every finite family $\mathcal F\subset A_+$ of bounded cylinder observables supported in $\{x_0>0\}$, the Wilson expectations
\[
\{\omega_{a_n}(F)\}_{F\in\mathcal F}
\]
converge jointly as $n\to\infty$. In particular, for every $F\in A_+$ the limit $\lim_{n\to\infty}\omega_{a_n}(F)$ exists and defines a state on $A_+$.
\end{corollary}

\begin{corollary}[Source Corollary~4.95: tuned dyadic convergence of reflected matrix elements]
Fix the same tuned dyadic trajectory and define the dense dyadic translation set
\[
D_{\mathrm{dyad}}:=\{m\ell_0 2^{-k}: m,k\in\mathbb N_0\}.
\]
Then for every finite family built from finitely many elements of $A_+$ by applying finitely many reflections and finitely many translations $\tau_s$ with $s\in D_{\mathrm{dyad}}$, the Wilson expectations converge jointly along the tuned dyadic sequence. In particular, for every bounded $F,G\in A_+$ and every $s\in D_{\mathrm{dyad}}$, the limit
\[
\lim_{n\to\infty}\omega_{a_n}(\Theta F\,\tau_s G)
\]
exists.
\end{corollary}

\begin{remark}
For the purposes of the core paper, the theorem-facing outputs of the present appendix are source Theorem~4.92 and Corollaries~4.93--4.95. Corollary~4.93 records the executive flowed-state formulation on $\mathcal A_{\mathrm{flow}}$, Corollary~4.94 gives the tuned bounded positive-time base state, and Corollary~4.95 gives the dyadic reflected matrix-element limits later consumed at the QE3 seam. The Wilson/RG identification on $\mathcal A_{\mathrm{flow}}$ is paired in the live suite with the flowed de-patching seam proved in Companion~I, Section~III. The larger dyadic-RG and injective-tuning machinery is retained here because it is part of the same native source block and because later Lane~A closure theorems continue to consume its outputs.
\end{remark}
